%==============================================================================%
% Chapter 5 – Field and Ring Formalisation
%==============================================================================%

\chapter{Formalisation of Field and Ring Arithmetic}
\label{ch:field_ring}

\section{The Prime Field $\F_q$}

\subsection{Field Definition in \texttt{Fq.ec}}

The file \texttt{Fq.ec} (536 lines) formalises the prime field
$\F_q = \Z / 64513\Z$.  Rather than working with \easycrypt{}'s built-in
\ec{ZModField} functor directly, the development uses a bespoke
\ec{clone} instantiation to gain tighter control over signed arithmetic.

\begin{lstlisting}[language=EasyCrypt,
  caption={Field declaration in \texttt{Fq.ec}},
  label={lst:fq_decl}]
require import AllCore IntDiv SmtMap Ring.
require import List StdBigop.

(* The modulus *)
op q : int = 64513.
axiom q_prime : prime q.
axiom q_pos   : 0 < q.

(* The canonical representative: representative in (-q/2, q/2] *)
op Fq_repr (a : int) : int = a %% q.
\end{lstlisting}

\subsection{Signed vs.\ Unsigned Representation}

A key design decision in \texttt{Fq.ec} is working with \emph{signed} integer
representatives in $(-q/2, q/2]$ rather than the conventional non-negative
representatives $[0, q)$.  This matches the signed 32-bit word representation
used in the \jasmin{} implementation, where coefficients are stored as
\jasm{u32} (unsigned 32-bit) but interpreted as signed integers via
two's-complement.

The file defines the \emph{signed word conversion} operator:
\[
  \op{to\_sint}(w) = \begin{cases}
    w & \text{if } w < 2^{31} \\
    w - 2^{32} & \text{otherwise}
  \end{cases}
\]
which maps a 32-bit unsigned word to its signed interpretation, and the
inverse \op{of\_sint} which lifts a bounded integer to a 32-bit word.

\subsection{Bit-Width Bounds}

A central technical device in \texttt{Fq.ec} is the \emph{bit-width bound
predicate}:
\[
  \op{bw32}(w, s) \;\Leftrightarrow\; \abs{\op{to\_sint}(w)} < 2^s,
\]
for a bit-width parameter $s \in \{1, \dots, 31\}$.  This captures the
invariant that a machine word represents a ``small'' integer.

\begin{lemma}[Bound addition]
  If $\op{bw32}(a, s_a)$ and $\op{bw32}(b, s_b)$, then
  $\op{bw32}(a + b, \max(s_a, s_b) + 1)$.
\label{lem:bw_add}
\end{lemma}

\begin{lemma}[Bound multiplication]
  If $\op{bw32}(a, s_a)$ and $\op{bw32}(b, s_b)$, then
  $\op{bw64}(a \cdot b, s_a + s_b)$
  where $\op{bw64}$ is the 64-bit analogue.
\label{lem:bw_mul}
\end{lemma}

These lemmas directly support the analysis of intermediate products in the
NTT butterfly: a 16-bit coefficient multiplied by a 15-bit twiddle factor
gives a 31-bit product, safe in a 32-bit signed word.

\subsection{Signed Arithmetic Lemmas}

The \easycrypt{} standard library does not include arithmetic lemmas for
signed integer interpretations of machine words.  \texttt{Fq.ec} contributes
several new lemmas, including:

\begin{lemma}[Signed right shift]
  For $a \in \Z$ with $\abs{a} < 2^{63}$,
  \[
    \op{to\_sint}(a \mathbin{\gg_s} 32) = a \div 2^{32},
  \]
  where $\gg_s$ denotes arithmetic right shift and $\div$ is truncating
  (towards zero) integer division.
\label{lem:arith_shift}
\end{lemma}

This lemma is crucial for formalising the Montgomery reduction formula, which
uses an arithmetic right shift to divide by $R = 2^{32}$.

\section{The \texorpdfstring{$\Z_q$}{Zq} Field via Theory Cloning}

\subsection{\texttt{GFq.ec}: Abstract Field}

The file \texttt{GFq.ec} clones \easycrypt{}'s abstract \ec{ZModField} theory
with $q = 64513$:

\begin{lstlisting}[language=EasyCrypt,
  caption={\texttt{GFq.ec}: Galois field instantiation},
  label={lst:gfq}]
require import ZModField IntDiv.
require Fq.

clone ZModField as Zq with
  op p <- Fq.q
  proof prime_p by exact Fq.q_prime.

(* Canonical embedding from int to Zq *)
abbrev inZq (a : int) : Zq.zmod = Zq.inzmod a.

(* Signed representative *)
op toInt (x : Zq.zmod) : int = Zq.asint x.
\end{lstlisting}

This instantiation provides the complete algebraic structure of $\F_q$:
commutativity, associativity, distributivity, existence of inverses, and the
characteristic-$q$ axiom $q \cdot 1 = 0$.

\subsection{Primitive Root}

The primitive $512$th root of unity $\zeta_0 = 426$ is defined and verified:

\begin{lstlisting}[language=EasyCrypt,
  caption={Primitive root declaration in NTTFullAlgebra.ec},
  label={lst:zroot}]
op zroot : Zq.zmod = inZq 426.

lemma zroot_512 : zroot ^ 512 = Zq.one.
lemma zroot_256 : zroot ^ 256 = Zq.(-\!one).
lemma zroot_primitive : forall k, 0 < k < 512 =>
    zroot ^ k <> Zq.one.
\end{lstlisting}

These are proved by direct computation in $\F_{64513}$, using the SMT
backend to discharge the modular arithmetic obligations.

\section{The Polynomial Ring $R_q$}

\subsection{\texttt{Rq.ec}: Ring Definition}

The file \texttt{Rq.ec} (118 lines) defines the polynomial ring
$R_q = \Z_q[X]/(X^{256}+1)$ as a 256-element coefficient vector over $\Z_q$:

\begin{lstlisting}[language=EasyCrypt,
  caption={Ring definition in \texttt{Rq.ec}},
  label={lst:rq_def}]
require import Array256.

(* Polynomial = coefficient vector *)
type poly = Zq.zmod Array256.t.

(* Ring addition *)
op (&+) (p q : poly) : poly =
  Array256.map2 Zq.(+) p q.

(* Ring multiplication (schoolbook, negacyclic) *)
op (&*) (p q : poly) : poly =
  let r = Array256.create Zq.zero in
  let r = foldr (fun i r =>
    let pi = p.[i] in
    foldr (fun j r =>
      let pij = pi * q.[j] in
      let idx = (i + j) %% 512 in
      if idx < 256 then r.[idx <- r.[idx] + pij]
                   else r.[idx - 256 <- r.[idx - 256] - pij]
    ) r (iota_ 0 256)
  ) r (iota_ 0 256) in
  r.
\end{lstlisting}

\subsection{NTT Axioms in \texttt{Rq.ec}}

Rather than proving from first principles that the NTT is a ring homomorphism,
\texttt{Rq.ec} introduces this as an \emph{axiom} that is later discharged
by the algebraic development in \texttt{NTTFullAlgebra.ec}:

\begin{lstlisting}[language=EasyCrypt,
  caption={NTT homomorphism axioms in \texttt{Rq.ec}},
  label={lst:rq_ntt_axiom}]
(* NTT of a product = pointwise product of NTTs *)
axiom ntt_mul_hom (f g : poly) :
  NTT (f &* g) = Array256.map2 Zq.( * ) (NTT f) (NTT g).

(* INTT is the left inverse of NTT *)
axiom intt_ntt (f : poly) : INTT (NTT f) = f.
\end{lstlisting}

These axioms are discharged as theorems once \texttt{NTTFullAlgebra.ec}
establishes the connection between the imperative NTT loop and the abstract
spectral map.

\section{Array Infrastructure}

\subsection{\texttt{Array256.ec}}

The 256-element array type \ec{Array256.t} is a typed abstract type with
the following operations:
\begin{itemize}
  \item \ec{.[i]} and \ec{.[i <- v]}: functional get and set.
  \item \ec{map}, \ec{map2}: pointwise operations.
  \item \ec{iteri}: iteration with index.
  \item \ec{init}: create from a function.
  \item \ec{create}: create with constant value.
\end{itemize}

Key properties include extensional equality (\ec{Array256.ext\_eq}):
two arrays are equal if and only if they agree at every index.

\subsection{Word Array Conversion}

The bridge between machine-word arrays and integer arrays is captured by
the following predicate:

\begin{definition}[Word-to-coefficient correspondence]
  Given a 32-bit word array $\mathtt{rp} : \ec{W32.t}^{256}$ and an integer
  array $\mathtt{a} : \Z^{256}$, we write $\mathtt{rp} \approx \mathtt{a}$
  (notation \ec{warray\_rel rp a}) if:
  \[
    \forall i \in [0, 256).\;
    \op{W32.to\_sint}(\mathtt{rp}[i]) = \mathtt{a}[i].
  \]
\label{def:wrel}
\end{definition}

This relation is the key linking the extracted \jasmin{} module (which works
with \ec{W32.t} arrays) to the imperative specification (which works with
\ec{int} arrays).  It is the core of the refinement proof in
\texttt{RefJasminNTT.ec}.

\section{Reduction and Auxiliary Operations}

\subsection{Conditional Addition (\texttt{caddq})}

The \ec{caddq} operation adds $q$ to negative coefficients:
\[
  \op{caddq}(a) = \begin{cases} a + q & \text{if } a < 0 \\ a & \text{otherwise} \end{cases}
\]
It maps signed representatives in $(-q, 0) \cup [0, q)$ to non-negative
representatives in $[0, q)$, without modular reduction.

\begin{lemma}[caddq correctness]
  For all $a \in \Z$ with $-q < a < q$:
  $\op{caddq}(a) \equiv a \pmod{q}$ and $0 \le \op{caddq}(a) < q$.
\label{lem:caddq}
\end{lemma}

\subsection{Freeze (\texttt{freeze})}

Full modular reduction to $[0, q)$ is defined as:
\[
  \op{freeze}(a) = ((a \bmod q) + q) \bmod q.
\]
\begin{lemma}[freeze correctness]
  For all $a \in \Z$:
  $\op{freeze}(a) \equiv a \pmod{q}$ and $0 \le \op{freeze}(a) < q$.
\label{lem:freeze}
\end{lemma}
